\section{Evaluation}
\label{sec:evaluation}

\subsection{Experimental Setup}

All experiments use the dataset constructed in Section~\ref{sec:dataset},
with a 70/10/20 stratified train/validation/test split and a separate
out-of-distribution (OOD) hold-out comprising the \texttt{empire\_mimikatz\_extract\_keys}
attack family, withheld entirely from training. This OOD set tests whether the
model generalises to unseen credential-access attack patterns beyond those seen
during fine-tuning.

The primary evaluation metric is the \textbf{False Negative Rate (FNR)} ---
the fraction of true attacks that the system fails to detect. In the LOLBin
detection domain, a missed attack is categorically more costly than a false
alarm, making FNR the correct objective metric. Secondary metrics are
Recall, Precision, F1-score, and AUC-ROC. All thresholds for the RF Tier~1
model are tuned on the validation set to achieve $\geq$98\% recall before
evaluating on the test set; all deep model thresholds are fixed at 0.5.

All experiments are implemented in PyTorch~2.x and scikit-learn~1.3,
with DistilBERT loaded from Hugging Face \texttt{distilbert-base-uncased}.
Training used a single NVIDIA GPU. Code and datasets will be released upon
publication.

\subsection{Comparison with Baseline Systems}
\label{sec:baselines}

Table~\ref{tab:comparison} compares our proposed hierarchical system against
four baselines evaluated on the same held-out test set.

\textbf{Rule-based (B1)} simulates the detection logic of LOTLDetector~\cite{zhu2026lotldetector},
applying a keyword and pattern-matching rule set derived from the LOLBAS project
without any learned model. \textbf{RF only (B2)} uses our Tier~1 Random Forest
as a standalone detector at its tuned threshold, representing conventional
ML-based LOLBin detection. \textbf{NLP only (B3)} applies our fine-tuned
DistilBERT model without graph context, representing the approach of Hendler
et al.~\cite{hendler2020amsi} and Yang et al.~\cite{yang2023robust}.
\textbf{GAT only (B4)} uses only the graph attention network with no command-line
NLP, representing the approach of Choi~\cite{choi2021powershell}.

Our full system achieves a recall of 0.9655, a precision of 0.9333,
an F1-score of 0.9492, and an AUC-ROC of 0.9871, with a false negative
rate of 0.034483 and 2 missed attacks on the test set.
By comparison, the rule-based baseline achieves FNR~=~0.275862,
the NLP-only baseline achieves FNR~=~0.103448, and the GAT-only
baseline achieves FNR~=~0.086207. These results demonstrate that
the full hierarchical architecture substantially reduces missed attacks
relative to any single-modality approach.

% Table 2 inserted here
\input{evaluation/baseline\_comparison\_table}

\subsection{Ablation Study}
\label{sec:ablation}

To isolate the contribution of each system component, we evaluate seven
configurations of the pipeline in Table~\ref{tab:ablation}.
Configurations~1--3 test each tier in isolation; Configurations~4--6
test pairwise combinations using equal-weight average ensembling;
Configuration~7 is the full system with the trained fusion MLP.

The most striking finding is the comparison between Config~6
(NLP~+~GAT without triage, FNR~=~0.068966) and Config~7
(full system, FNR~=~0.034483). Adding the RF triage layer does not
merely improve precision --- it reduces FNR by providing an initial
high-recall sweep that ensures no attack is discarded before reaching
the deep models. Removing the triage layer causes a disproportionate
increase in FNR because some attack sessions with short or unobfuscated
commands are missed by the NLP and GAT models when not first flagged
by the RF.

The comparison between Config~4 (RF~+~NLP) and Config~5 (RF~+~GAT) is
also instructive: RF~+~GAT achieves lower FNR than RF~+~NLP, confirming
that graph-based process-tree context is a stronger detection signal than
command-line content alone for the structural attack patterns targeted
in this work. However, the full system outperforms both, confirming that
content and context are complementary rather than redundant.

% Table 3 inserted here
\input{evaluation/ablation\_table}

\subsection{Out-of-Distribution Generalisation}
\label{sec:ood}

Table~\ref{tab:ood} compares IID and OOD performance across all model tiers.
The OOD set comprises credential-access attack families (\texttt{Mimikatz}-based
exfiltration chains) withheld entirely during training. A robust system should
show small $\Delta$FNR between IID and OOD evaluation.

Our full system achieves $\Delta$FNR~=~+0.0155 between the IID test set
and the OOD set, indicating strong generalisation. This robustness stems from
two architectural decisions: (1) the GAT component learns structural process-tree
patterns (parent-child spawn chains, LOLBin sequences) that generalise across
attack families, not family-specific command strings; and (2) the DistilBERT
component detects obfuscation patterns (Base64 encoding, execution policy bypass,
encoded payloads) that are present across attack families regardless of the
specific tool used. By contrast, the rule-based baseline (B1) shows the largest
$\Delta$FNR on OOD data, as expected for a system that relies on pattern
matching against known signatures.

% Table 4 inserted here
\input{evaluation/ood\_comparison\_table}

\subsection{Per-Family Attack Analysis}
\label{sec:family}

Figure~\ref{fig:family} shows FNR broken down by attack family.
The lowest FNR is observed for \texttt{empire\_psinject} (FNR~=~0.020000),
which produces highly characteristic process trees readily detected by the GAT.
The highest FNR is observed for \texttt{splunk\_T1003.001\_windows-sysmon} (FNR~=~0.071429),
where the attack chain is shorter and involves fewer LOLBin processes, reducing
the graph-structural signal available to the GAT. This finding motivates future
work on improving detection of short, low-complexity attack chains.

% Figure: family FNR heatmap
\begin{figure}[h]
\centering
\includegraphics[width=0.95\linewidth]{evaluation/family_fnr_heatmap}
\caption{FNR by attack family (full system). Red bars indicate families
where FNR exceeds 5\%.}
\label{fig:family}
\end{figure}

\subsection{Attack Chain Reconstruction Quality}
\label{sec:reconstruction}

We evaluate the quality of the automated MITRE ATT\&CK kill chain
reconstruction on all test-set sessions classified as malicious
by the fusion model. Reconstruction quality is measured by TTP Precision
(fraction of predicted TTPs matching ground-truth labels),
TTP Recall (fraction of ground-truth TTPs recovered),
and Chain Completeness (fraction of sessions with $\geq$2 distinct TTPs
in the reconstructed chain).

Results are reported in Table~\ref{tab:reconstruction} and representative
case-study chains are shown in Section~\ref{sec:casestudies}. The system
achieves a chain completeness of over 90\% on test-set malicious sessions,
meaning that for the vast majority of detected attacks, the system produces
an interpretable multi-step kill chain rather than a single-label alert.
This output is the primary differentiator of our system from prior detection
work: rather than issuing a binary alert, the system automatically reconstructs
the full attack narrative for analyst consumption.